\section{Background and Related Work}
\label{sec:relatedworkMain}

\subsection{Definition of Code Comprehensibility}


\subsection{Comprehension Tasks in Prior Work}
\label{sec:rw-tasks}

\begin{table}
\centering
\caption{Taxonomy of comprehension tasks based on Wyrich~\etal~\cite{Wyrich:ACS23}, with updated study counts through March 2026.}
\resizebox{\columnwidth}{!}{
\begin{tabular}{p{3cm} l c}
  \toprule
 Task Category & Task & Number of Studies \\ 
  \midrule

\multirow{6}{3cm}{Tc1: Provide information about the code}
    & Answer comprehension questions & 37 \\ 
    & Determine output of a program & 30 \\ 
    & Summarize code verbally or textually & 25 \\ 
    & Recall (parts of) the code & 7 \\ 
    & Match with diagram or similar code & 3 \\
    & Determine code trace & 1 \\ 
  \midrule

\multirow{6}{3cm}{Tc2: Provide personal opinion}
  & Rate code comprehensibility & 13 \\ 
  & Other subjective indications & 7 \\ 
  & Rate task difficulty or task load & 5 \\ 
  & Rate confidence in answer/understanding & 5 \\ 
  & Compare or rank code snippets & 5 \\
  & Rate own code understanding & 2 \\  
  \midrule

\multirow{2}{3cm}{Tc3: Debug code}
  & Find a bug & 9 \\ 
  & Extend the code & 4 \\ 
  \midrule

\multirow{4}{3cm}{Tc4: Maintain code}
    & Fix a bug & 4 \\ 
    & Cloze test on code & 3 \\ 
    & Refactor code & 3 \\ 
    & Determine if code is correct & 1 \\
 
   \bottomrule
\end{tabular}
}
\label{tab:proxies-wyrich}
\end{table}
 
\Cref{tab:proxies-wyrich} extends the taxonomy of comprehension tasks in Wyrich~\etal's survey~\cite{Wyrich:ACS23}
with papers published on or before March 25, 2026 but after their literature search.
While there is significant diversity in tasks, 
the literature has a strong bias toward \textbf{Tc1} and, to a lesser extent, 
\textbf{Tc2}. The most common tasks in these categories are answering questions of 
various types, free-text summarization, and subjective ratings; together, these 
tasks dominate the literature. 


\subsection{Comprehension Measures in Prior Work}
\label{sec:rw-measures}

\begin{table}[t]
\centering
\small
\renewcommand{\arraystretch}{1.1}
\caption{Frequency of comprehension measure types based on the taxonomy of Wyrich~\etal~\cite{Wyrich:ACS23}, with updated study counts through March 2026.}
\begin{tabular}{lclc}
\toprule
Measure Type & Studies & Measure Type & Studies \\
\midrule
Correctness & 80 & Physiological & 30 \\
Time & 50 & Aggregation & 6 \\
Subjective rating & 37 & Others & 6 \\
\bottomrule
\end{tabular}
\label{tab:measures}
\end{table}
 
\Cref{tab:measures} extends the comprehension measure survey data from Wyrich~\etal~\cite{Wyrich:ACS23} with
that from studies published before March 25, 2026.
\emph{Correctness} (also referred to as \emph{accuracy} ~\cite{S29}) compares a study participant's (\ie a subject) answer to a known ground truth~\cite{Nielebock:ESE2019, Peitek:ICPC20, Peitek:ICSE2021, Ribeiro:ESE23, Karas:TSE2024, Park:ESE24, Alakmeh:ASE2024, Elfares:2025, Flint:2026, Abdelsalam:ESE2026}.
\emph{Time} captures how long a subject takes to complete 
a task, regardless of answer correctness~\cite{Nielebock:ESE2019, Peitek:ICPC20, Peitek:ICSE2021, Karas:TSE2024, Park:ESE24, Alakmeh:ASE2024, Elfares:2025, Abdelsalam:ESE2026}.
\emph{Subjective ratings} are self-reported perceptions, such as 
Likert-scale responses to questions~\cite{Nielebock:ESE2019, Peitek:ICSE2021, Ribeiro:ESE23, Abdelsalam:ESE2026, Alakmeh:ASE2024, Elfares:2025, Bergum:TOSEM2026}.
\emph{Physiological} measurements capture responses of the human body or brain when performing code comprehensibility tasks. 
Eye-tracking~\cite{Peitek:ICPC20, Abbad-Andaloussi:ICPC22, Park:ESE24, Abdelsalam:ESE2026}, fMRI~\cite{Peitek:ICSE2021} and EEG~\cite{Bergum:TOSEM2026} are the most common measures in this category.
\emph{Aggregation} combines two or more other measures~\cite{Ribeiro:ESE23, Alakmeh:ASE2024, Gao:EASE2025}; 
for example, ``time to correctly answer a question'' aggregates correctness 
and time measures.
\emph{Other} includes measures that do not fit into the above categories, such as using code metrics to measure comprehensibility~\cite{33-8, Nielebock:ESE2019}.



